% GENERATED from content/data/enrichment.json by mxbuild.enrichment; do not edit.
\mxunit{enrichment-bisection-exercises}{Encadrer une racine par dichotomie}
\mxsubtitle{Terminale · analyse et algorithmique \textperiodcentered{} Exercices progressifs \textperiodcentered{} 12}

\begin{note}
\textbf{Données.} Pour l’exercice 1 : on cherche un zéro de $f(x)=x^{2}-2$ sur $\intff{1}{2}$. Après deux partages, l’intervalle conservé est $\intff{\num{1,25}}{\num{1,5}}$.
\end{note}

\begin{exercise}
Poursuivre l’exemple une troisième fois.
\end{exercise}

\begin{exercise}
Partant d’un intervalle de largeur 1, combien de partages garantissent une largeur au plus $\frac{1}{1000}$ ?
\end{exercise}

\begin{exercise}
Pourquoi le changement de signe de $f(x)=\frac{1}{x}$ entre $-1$ et 1 ne prouve-t-il pas l’existence d’un zéro ?
\end{exercise}

\begin{exercise}
Une racine est dans $\intff{\num{1,4}}{\num{1,5}}$. Quelle erreur maximale garantit le milieu $\num{1,45}$ ?
\end{exercise}

